\begin{titlepage}
\thispagestyle{empty}
\begin{center}
{\huge\bfseries First Steps\\[0.2cm]}
{\Large Empirical Study of Public Law Litigation\\[0.35cm]}
{\large\itshape in Sri Lanka's Apex Courts\\[1.2cm]}

{\large Elijah Hoole\textsuperscript{$\dagger$}, Dinesha Samararatne\textsuperscript{$\ddagger$} \& Yanitra Kumaraguru\textsuperscript{$\S$}\\[0.8cm]}

{\small\scshape Dataset}\\[0.25cm]
{\large\texttt{SriLankaCourtDecisions} (2012--2022)\\[1cm]}

{\small Draft generated on \today}
\end{center}
\vfill

\noindent\rule{\textwidth}{0.4pt}\\[0.3cm]
{\footnotesize
\textsuperscript{$\dagger$} Founder, paralegal.lk (lead author, contributed approximately 50\% towards this article)\\[0.15cm]
\textsuperscript{$\ddagger$} Professor, Dept of Public \& International Law, Faculty of Law, University of Colombo, Sri Lanka and independent member of the Constitutional Council of Sri Lanka (Jan 2023 -- Jan 2025) (contributed approximately 40\% towards this article)\\[0.15cm]
\textsuperscript{$\S$} Team Lead, Legal Research, Verité Research (contributed approximately 10\% towards this article)\\[0.3cm]
The research work for this paper was carried out 2023 to 2025. A preliminary draft of this paper by Elijah Hoole, Amanda Halliday \& Yanitra Kumaraguru `Empirical Study of Sri Lanka's Apex Courts: Taking the First Step with Dataset Curation and Analysis' was presented by Amanda Halliday at the conference on Computational Analysis of Apex Courts 2023 held June 23, in Braga, Portugal. We thank Tavini Nanayakkara, Akalanka Tilakeratne and Shanil Wijesinha for their feedback.
}
\end{titlepage}
\pagenumbering{roman}
\setcounter{page}{1}
